\documentclass[11pt]{article}

\usepackage[margin=1in]{geometry}
\usepackage[T1]{fontenc}
\usepackage[utf8]{inputenc}
\usepackage{lmodern}
\usepackage{microtype}
\usepackage{graphicx}
\usepackage{booktabs}
\usepackage{enumitem}
\usepackage{hyperref}
\usepackage{tikz}
\usetikzlibrary{positioning,arrows.meta}

\title{Architecture Writeup\\[0.3em]
\large DSCI 601 -- Applied Data Science I}
\author{Piter Z. Garcia}
\date{April 18, 2026}

\begin{document}
\maketitle

\section{Overview}
This architecture writeup describes the implementation of the code review repository
used for the DSCI 601 assignment. The repository is organized as a small,
reproducible \emph{Hybrid CRISP+ML} pipeline that mirrors a real machine-learning
workflow while remaining fast enough for reviewers to install, run, test, and
inspect locally. The main design goal was to create a reviewer-safe system with
clear stage boundaries, deterministic behavior, and enough implementation detail
that another party could recreate the workflow and reproduce the same outputs.

The system is structured around seven explicit lifecycle stages:
\begin{enumerate}[leftmargin=1.5em]
    \item Data Collection
    \item Data Cleaning
    \item Data Preparation
    \item Data Exploration / Visualization
    \item Data Mining
    \item Evaluation
    \item Results Postprocessing / Visualization
\end{enumerate}
These stages were chosen so the implementation matches both a recognized
methodology and the assignment requirements, which ask for separable runnable code
for data handling, modeling, exploration, and postprocessing.

\section{Architecture and Layer Interactions}
Figure~\ref{fig:architecture} shows the overall system structure. The software is
organized into four main layers: configuration, stage execution, orchestration,
and outputs.

\begin{figure}[h]
\centering
\resizebox{\textwidth}{!}{% Architecture diagram for the DSCI601 architecture report.
% The figure mirrors the implemented CRISP+ML stage pipeline and highlights
% how configuration, CLI commands, tests, and storage wrap the stages.

\begin{tikzpicture}[
  font=\sffamily,
  >=Latex,
  support/.style={
    rectangle,
    draw=black,
    rounded corners,
    minimum height=0.9cm,
    text width=3.6cm,
    align=center,
    fill=gray!10,
    line width=0.45pt
  },
  stage/.style={
    rectangle,
    draw=black,
    rounded corners,
    minimum height=1.25cm,
    text width=3.0cm,
    align=center,
    line width=0.45pt
  },
  store/.style={
    rectangle,
    draw=black,
    rounded corners,
    minimum height=0.9cm,
    text width=4.6cm,
    align=center,
    fill=gray!8,
    line width=0.45pt
  },
  arrow/.style={->, line width=0.75pt},
  dashedarrow/.style={->, dashed, line width=0.6pt},
  note/.style={font=\scriptsize, align=center}
]

% Support / control layer
\node[support, fill=phase1blue!12] (config) at (0, 2.1)
  {\textbf{Configuration}\\seed, split, missingness,\\fairness threshold};
\node[support, fill=phase2green!12] (cli) at (5.0, 2.1)
  {\textbf{CLI Interface}\\one command per stage\\\texttt{run-all}, \texttt{merge-check}};
\node[support, fill=phase4purple!12] (tests) at (10.0, 2.1)
  {\textbf{Validation Layer}\\unit tests, smoke run,\\CI merge gate};

% CRISP+ML lifecycle row 1
\node[stage, fill=phase1blue!15] (collect) at (0, 0)
  {\textbf{1. Data}\\\textbf{Collection}\\synthetic reviewer dataset\\\scriptsize output: raw JSON};
\node[stage, fill=phase1blue!15] (clean) at (3.8, 0)
  {\textbf{2. Data}\\\textbf{Cleaning}\\schema validation\\\scriptsize output: clean JSON};
\node[stage, fill=phase2green!15] (prepare) at (7.6, 0)
  {\textbf{3. Data}\\\textbf{Preparation}\\deterministic split\\\scriptsize output: train/test JSON};
\node[stage, fill=phase2green!15] (explore) at (11.4, 0)
  {\textbf{4. Exploration}\\\textbf{/ Visualization}\\group + completeness\\summaries\\\scriptsize output: CSV};

% CRISP+ML lifecycle row 2
\node[stage, fill=phase3orange!15] (mine) at (0, -2.3)
  {\textbf{5. Data}\\\textbf{Mining}\\lightweight model run\\\scriptsize output: predictions JSON};
\node[stage, fill=phase3orange!15] (evaluate) at (3.8, -2.3)
  {\textbf{6. Evaluation}\\utility + fairness metrics\\\scriptsize output: metrics JSON};
\node[stage, fill=phase4purple!15] (post) at (7.4, -2.3)
  {\textbf{7. Results Postprocess}\\final tables + SVG\\\scriptsize output: CSV + figure};
\node[store] (manifest) at (11.5, -2.3)
  {\textbf{Pipeline Manifest}\\config copy, stage outputs,\\summary metadata};

% Stage flow
\draw[arrow] (collect) -- (clean);
\draw[arrow] (clean) -- (prepare);
\draw[arrow] (prepare) -- (explore);
\draw[arrow] (mine) -- (evaluate);
\draw[arrow] (evaluate) -- (post);
\draw[arrow] (post) -- (manifest);

% Control / validation flow
% Storage/provenance note
\node[note, align=center] at (5.7, -3.85) {
  \textbf{Storage contract:} each stage writes deterministic artifacts under
  \texttt{results/fast\_review/<numbered-stage>/}; the manifest records the full run.
};

\end{tikzpicture}
}
\caption{Architecture of the reviewer-safe CRISP+ML pipeline.}
\label{fig:architecture}
\end{figure}

\subsection{Configuration Layer}
The configuration layer defines the parameters that control the entire run. In the
implementation repository, this is handled by a typed configuration object loaded
from a single JSON file. The configuration stores the random seed, record counts,
train/test split, invalid-record frequency, missingness frequency, bias level,
fairness threshold, and output directory.

This layer was separated from the stage logic for two reasons. First, it keeps
experiments reproducible: the same configuration and seed lead to the same output
artifacts. Second, it makes the system easy to inspect and modify without changing
source code. A reviewer can understand the workflow by reading one config file
instead of tracing hard-coded constants across the repository.

\subsection{Stage Layer}
Each lifecycle step is implemented as its own concrete class under a shared
\texttt{PipelineStage} interface. This makes the code strongly object-oriented and
keeps stage responsibilities narrow. The stage implementations behave as follows:

\begin{itemize}[leftmargin=1.5em]
    \item \textbf{Data Collection} generates a small synthetic dataset that
    preserves fairness-relevant structure such as group membership, label, and
    feature availability.
    \item \textbf{Data Cleaning} validates records, removes intentionally corrupted
    examples, and standardizes the remaining schema.
    \item \textbf{Data Preparation} converts cleaned records into model-ready train
    and test artifacts.
    \item \textbf{Data Exploration} creates lightweight summaries of dataset size,
    group distribution, and feature completeness.
    \item \textbf{Data Mining} runs the minimal modeling step on the prepared
    dataset and writes predictions.
    \item \textbf{Evaluation} computes utility and fairness-oriented metrics from
    predictions.
    \item \textbf{Results Postprocessing} turns evaluation summaries into reviewer-
    friendly output tables and simple figures.
\end{itemize}

Every stage reads only the artifacts it needs, writes to its own dedicated output
directory, and returns a structured result object containing artifact paths and
summary metadata. This design keeps the stages testable in isolation and reduces
the risk of hidden side effects between steps.

\subsection{Context and Orchestration Layer}
The orchestration layer is responsible for connecting the individual stages
together. A shared pipeline context object carries metadata, output locations, and
stage results through the run so that stages do not need direct knowledge of each
other's internal implementation. This context acts as a lightweight contract for
information sharing.

The \texttt{ReviewPipeline} orchestrator builds the ordered list of stages and
executes them in sequence. It can run the full workflow or stop after any named
stage. At the end of execution, it writes a manifest file that records the active
configuration, stage outputs, and stage summaries. This manifest becomes the
single reproducibility record for a run.

\subsection{Interface Layer}
The user-facing interface is a command-line entry point with one command for each
stage plus two orchestration commands: one for running the full pipeline and one
for performing the local merge validation path. This layer was important for the
assignment because the repository needed explicit runnable scripts rather than a
single notebook. The command-line interface lets another party run individual
phases independently or validate the full workflow end-to-end.

\section{Framework and Component Choices}
The architecture intentionally uses a very small technical stack. The pipeline is
implemented in standard-library Python with a minimal package structure and
\texttt{requirements.txt} as the reviewer-facing install path. This choice was
made to reduce environmental friction during code review. Reviewers do not need to
install a heavy machine-learning stack or resolve hidden dependencies just to run
the repository.

JSON was selected for configuration and most structured artifacts because it is
easy to inspect, easy to validate, and natively supported in Python. CSV was used
for summary tables because it is a familiar exchange format for grading and quick
inspection. SVG was used for lightweight visual outputs because it provides a
clear figure artifact without requiring plotting libraries. The overall component
selection favors reproducibility, transparency, and short setup time over raw
performance.

\section{Data Storage and Output Design}
All artifacts are written beneath a stable output root defined by the active
configuration. Within that root, each stage writes to its own numbered directory.
This output scheme was chosen for traceability. A reviewer can inspect any stage
without rerunning the full system, and each run naturally leaves behind a visible
artifact trail.

The output structure is designed to answer three practical questions:
\begin{itemize}[leftmargin=1.5em]
    \item What input configuration produced this run?
    \item What did each stage write?
    \item What metrics and final outputs were generated?
\end{itemize}

Those questions are answered by the combination of stage directories and the final
manifest file. This layout also supports merge validation because the validation
step can check for expected artifact existence in deterministic locations.

\section{Reproducibility Path}
Another party can recreate the system by setting up a fresh Python environment,
installing the repository, and executing the pipeline with the provided fast
review configuration. The expected workflow is:
\begin{enumerate}[leftmargin=1.5em]
    \item install dependencies from \texttt{requirements.txt},
    \item run the full pipeline command with the provided configuration,
    \item inspect the generated results directory,
    \item run the merge-check command to execute the local validation path.
\end{enumerate}

This path reproduces the same stage artifacts, metrics, summaries, and validation
checks that the repository uses for review. Because the system is configuration-
driven, deterministic, and documented through explicit stage boundaries, another
party should be able to recreate both the software behavior and the presented
results with minimal ambiguity.

\section{Conclusion}
The architecture emphasizes simplicity, determinism, and inspectability. Rather
than optimizing for scale, the design optimizes for a clean code review experience:
small inputs, explicit stages, typed configuration, structured outputs, and a
repeatable validation path. This makes the system appropriate both for the DSCI
601 assignment and for future extension into larger reproducible workflows.

\end{document}
